% Submitted on ... to ...
% Revised on ...
\documentclass[a4paper]{article}
%\documentclass[a4paper,titlepage]{article}

\usepackage{amsmath}
\usepackage{amssymb}
%\usepackage[dvips]{graphicx}
\usepackage{graphicx}
\usepackage[OT4]{fontenc}
%\usepackage[T1]{fontenc}
\usepackage[utf8]{inputenc}
\usepackage{tikz}
\usepackage{url}
\usepackage[numbers]{natbib}
\usepackage{bm}
\emergencystretch=20pt

\usepackage{setspace}
%\doublespacings

\newcommand{\abs}[1]{\left| #1 \right|}
\newcommand{\norm}[1]{\left|\mkern-1mu\left| #1 \right|\mkern-1mu\right|}
\newcommand{\aabs}[1]{\left\| #1 \right\|}
\newcommand{\okra}[1]{\left( #1 \right)}
\newcommand{\kwad}[1]{\left[ #1 \right]}
\newcommand{\klam}[1]{\left\{ #1 \right\}}
\newcommand{\ceil}[1]{\left\lceil #1 \right\rceil}
\newcommand{\floor}[1]{\left\lfloor #1 \right\rfloor}
\newcommand{\dzi}[1]{\left\langle #1 \right\rangle}
\newcommand{\ind}{\mathrel{\perp\mkern-10mu\perp}}
\DeclareMathOperator*{\argmin}{arg\, min}
\DeclareMathOperator*{\argmax}{arg\, max}
\DeclareMathOperator*{\hilberg}{hilb}
\DeclareMathOperator{\essinf}{ess\, inf}
\DeclareMathOperator{\esssup}{ess\, sup}
\DeclareMathOperator{\card}{card}
\DeclareMathOperator{\supp}{supp}
\DeclareMathOperator{\cov}{Cov}
\DeclareMathOperator{\var}{Var}
\DeclareMathOperator{\mean}{\mathbf{E}}
\DeclareMathOperator{\median}{\mathbf{M}}
\newcommand{\peq}{\stackrel{+}{=}}
\newcommand{\pge}{\stackrel{+}{>}}
\newcommand{\ple}{\stackrel{+}{<}}
\newcommand{\boole}[1]{{\bf 1}{\klam{#1}}}

\newtheorem{conjecture}{Conjecture}
\newtheorem{definition}{Definition}
\newtheorem{example}{Example}
\newtheorem{corollary}{Corollary}
\newtheorem{proposition}{Proposition}
\newtheorem{theorem}{Theorem}
\newtheorem{lemma}{Lemma}
\newenvironment*{proof}{\begin{trivlist}\item[]
\noindent\textbf{Proof:}}{$\Box$\par\end{trivlist}}
\newenvironment*{proof*}[1]{\begin{trivlist}\item[]
\noindent\textbf{Proof of #1:}}{$\Box$\par\end{trivlist}}
%\renewcommand{\theenumi}{\roman{enumi}}
%\renewcommand{\labelenumi}{(\theenumi)}
%\renewcommand{\thesection}{\Roman{section}}
%\renewcommand{\thesubsection}{\Roman{section}-\Alph{subsection}}


\newcommand{\bibnamefont}[1]{\MakeUppercase{#1}}


\author{{\L}ukasz D\k{e}bowski\thanks{
    {\L}. D\k{e}bowski is with
    the Institute of Computer Science, Polish Academy of Sciences, 
    ul. Jana Kazimierza 5, 01-248 Warszawa, Poland 
    (e-mail: ldebowsk@ipipan.waw.pl).
  }
  % \thanks{Manuscript received ... ; revised ... .}
  %\thanks{}  
}


\title{Recurrence and repetition times in the case of \\ a stretched
  exponential growth} \date{}

\begin{document}

%\pagestyle{empty}   
%\begin{titlepage}
%\titlepage
\maketitle

\begin{abstract}
  By an analogy to the duality between the recurrence time and the
  longest match length, we introduce a quantity dual to the maximal
  repetition length, which we call the repetition time.  Using the
  generalized Kac lemma for successive recurrence times by Chen Moy,
  we sandwich the repetition time in terms of min-entropies with no or
  relatively short conditioning. The sole assumption is stationarity
  and ergodicity. The proof is surprisingly short and the claim is
  fully general in contrast to earlier approaches by Szpankowski and
  by D\k{e}bowski. We discuss the analogy of this result with the
  Wyner-Ziv/Ornstein-Weiss theorem, which sandwiches the recurrence
  time in terms of Shannon entropies. We formulate the respective
  sandwich bounds in a way that applies also to the case of stretched
  exponential growth observed empirically for natural language.
% \end{abstract}
%
% \begin{keywords}
  \\[1em]
  \textbf{Keywords:} stationary processes, recurrence time, maximal
  repetition, R\'enyi entropies
  % \end{keywords}
  \\[1em]
  \textbf{MSC2020:} 60G10, 94A17
\end{abstract}

% %MSC 2020: 
% % 94A29 - Source coding
% % 60G10 - Stationary processes
% % 94A17 - Measures of information, entropy


% \end{titlepage}
% \pagestyle{plain}   

% \tableofcontents

\section{Introduction}
\label{secIntroduction}

Using appropriate concepts is often a gate to the simple theory. Some
instance of this rule will be exhibited in this short note.  Let
$(X_i)_{i\in\mathbb{Z}}$ be an arbitrary stochastic process over a
countable alphabet. The well-researched recurrence time $R^{(1)}_k$
\citep{Kac47,WynerZiv89,OrnsteinWeiss93,Kontoyiannis98,Ko12} is the
first position in the future on which a copy of string $X_1^k$
reappears,
\begin{align}
  R^{(1)}_k:=\inf\klam{\infty}\cup\klam{i\ge 1: X_{i+1}^{i+k}=X_1^k}.
\end{align}
By contrast, we define here the repetition time $R^{(2)}_k$ as the
first position in the future on which a copy of any previously seen
string $X_{j+1}^{j+k}$ occurs,
\begin{align}
  R^{(2)}_k:=\inf\klam{\infty}\cup\klam{i\ge 1:
  X_{i+1}^{i+k}=X_{j+1}^{j+k}\text{ for some } 0\le j<i}.
\end{align}
If $(X_i)_{i\in\mathbb{Z}}$ is a process over a $D$-ary alphabet then
$R^{(2)}_k\le D^k$. In any case, $R^{(2)}_k\le R^{(1)}_k$. Precisely,
we have
\begin{align}
  \label{Minimum}
  R^{(2)}_k=\inf_{i\ge 0}\okra{i+R^{(1)}_k\circ T^i}
\end{align}
for the shift automorphism $X_i\circ T=X_{i+1}$, since observing a
repetition is equivalent to observing a recurrence of some particular
string. Relationship (\ref{Minimum}) is a simple fact that will have
weighty consequences in this note.
% The motivation for the novel term ``impossibility time'' is a rule
% of thumb that if some rare event happens twice then the phenomenon
% cannot be random.

In parallel, we can discuss variables that are dual to random
variables $R^{(\gamma)}_k$. The longest match length $L^{(1)}_n$
\citep{WynerZiv89,OrnsteinWeiss93,Szpankowski93,
  KontoyiannisOthers98,GaoKontoyiannisBienenstock08}
is
\begin{align}
  L^{(1)}_n:=\max\klam{0}\cup\klam{k\ge 1: X_{i+1}^{i+k}=X_1^k
  \text{ for some } 0< i\le n-k},
\end{align}
whereas the maximal repetition length $L^{(2)}_n$
\citep{DeLuca99,KolpakovKucherov99a,KolpakovKucherov99,CrochemoreIlie08,
  ErdosRenyi70,ArratiaWaterman89,Shields92b,Shields97,
  Szpankowski93,Szpankowski93a,KontoyiannisSuhov94,
  Debowski11b,Debowski17,Debowski18b} is
\begin{align}
  L^{(2)}_n:=\max\klam{0}\cup\klam{k\ge 1: X_{i+1}^{i+k}=X_{j+1}^{j+k}
  \text{ for some } 0\le j< i\le n-k}.
\end{align}
Both $L^{(\gamma)}_n$ are functions of prefix $X_1^n$ and
$n\ge L^{(2)}_n\ge L^{(1)}_n$. For $\gamma\in\klam{1,2}$, we have the
dualities
\begin{align}
  R^{(\gamma)}_k&:=\min\klam{n: L^{(\gamma)}_{n+k}\ge k},
  &
    L^{(\gamma)}_n&=\max\klam{k: R^{(\gamma)}_k\le n-k}.
\end{align}
Hence $R^{(\gamma)}_k>n\iff L^{(\gamma)}_{n+k}<k$
\citep{WynerZiv89}. Thus the dual random variables carry the same
information as the primal. Both $R^{(\gamma)}_k$ and $L^{(\gamma)}_n$
are non-decreasing functions of $k$ and $n$ respectively.

Whereas variable $R^{(2)}_k$ is defined in this paper, there is a long
history of investigation of variables $R^{(1)}_k$, $L^{(1)}_n$, and
$L^{(2)}_n$. These concepts appear in various contexts in ergodic
theory, information theory, theoretical computer science, and
quantitative linguistics, as we have already indicated by the
references.  Establishing the asymptotic behavior of $R^{(\gamma)}_k$
and $L^{(\gamma)}_n$ for stationary ergodic processes is of a theoretical
importance since combined with empirical data, it can reveal what kind
of a process may generate these observations.

To comment on the respective theoretical results, let us begin with
the recurrence time and the longest match length.  Let $\log$ denote
the natural logarithm.  Consider a discrete random variable $X$. The
Shannon entropy is $H_1(X):=\mean\kwad{-\log P(X)}$. Let
$(X_i)_{i\in\mathbb{Z}}$ be a stationary ergodic process over a finite
alphabet. The celebrated Wyner-Ziv/Ornstein-Weiss (WZ/OW) theorem
\citep{WynerZiv89,OrnsteinWeiss93} asserts
\begin{align}
  \label{WZOW}
  \lim_{k\to\infty} \frac{\log n}{L^{(1)}_n}
  =\lim_{k\to\infty} \frac{1}{k}\log R^{(1)}_k
  =h_1:=\lim_{k\to\infty} \frac{H_1(X_1^k)}{k}
  \text{ a.s.}
\end{align}
Thus the rates of the recurrence time and the longest match length are
essentially dictated by the rate of the Shannon entropy.  The Shannon
entropy rate $h_1$ exists for any stationary process.

Let us proceed to the maximal repetition length and the repetition
time.  For memoryless sources, the maximal repetition length
$L^{(2)}_n$ grows at a similar rate as the longest match length
$L^{(1)}_n$ \citep{ErdosRenyi70,ArratiaWaterman89}. Probably for this
reason, the problem that $L^{(2)}_n$ can grow much faster than
$L^{(1)}_n$ for other stationary processes was first recognized as
late as by \citet{Shields92b,Shields97}. Define the collision entropy
is $H_2(X)=-\log\mean P(X)$, where $H_2(X)\le H_1(X)$ by the Jensen
inequality.  A technically complicated result by
\citet{Szpankowski93,Szpankowski93a}, involving suffix trees, can be
rewritten as
\begin{align}
  \label{Szpankowski}
  \lim_{k\to\infty} \frac{\log n}{L^{(2)}_n}
  =\lim_{k\to\infty} \frac{1}{k}\log R^{(2)}_k
  =h_2:=\lim_{k\to\infty} \frac{H_2(X_1^k)}{k}
  \text{ a.s.}
\end{align}
provided we have a uniform bound on the pointwise mutual information
\begin{align}
  \label{Kfinite}
  \sup_{n,m\ge 1}\esssup
     \frac{P(X_1^{n+m})}{P(X_1^n)P(X_{n+1}^{n+m})}
     < \infty.
\end{align}
Quantity $h_2$ is called the collision entropy rate. Whereas limit
$h_1$ exists for any stationary process, limit $h_2$ exists if
(\ref{Kfinite}) holds additionally. The collision entropy rate can be
substantially smaller than the Shannon entropy rate.

Condition (\ref{Kfinite}) is satisfied by finite-state processes
\citep[Theorem 11.3]{Debowski21} but it does not cover all important
cases---such as stationary sources that do not satisfy the finite
energy condition \citep{Shields92b,Shields97}, do not meet the
Doeblin condition \citep[Theorem 11.7]{Debowski21}, or are not diluted
by random noise \citep{Shields97}, \citep[Theorem 11.6]{Debowski21}.
We became interested in characterizing the maximal repetition length
because of our long-standing interests in statistical language
modeling \citep{Debowski21}.  In fact, for natural language,
experiments suggest that the Shannon entropy rate is positive,
$h_1>0$, \citep{Shannon51,CoverKing78,BehrOthers03} but the mutual
information probably grows like a power-law, called Hilberg's law
\citep{Hilberg90, TakahiraOthers16, KaplanOther20, Debowski21}, which
contradicts condition (\ref{Kfinite}). Moreover, as discovered
empirically in \citep{Debowski12b,Debowski15f} for novels in English,
German, and French, we have a power-law logarithmic growth of the
maximal repetition length
\begin{align}
  L^{(2)}_n\propto \log^\alpha n,\quad \alpha\approx 3.
\end{align}
This is equivalent to the stretched exponential growth of the
repetition time
\begin{align}
  \log R^{(2)}_k\propto k^\beta,\quad \beta\approx 1/3.
\end{align}

The goal of this paper is to analyze this case of a stretched
exponential growth, rectifying and beautifying our earlier theoretical
accounts and cumbersome statements
\citep{Debowski15f,Debowski18b,Debowski21}. We will do it by using the
repetition time $R^{(2)}_k$ rather than the maximal repetition length
$L^{(2)}_n$. The new results are simple and general but they have
escaped our awareness for quite a long time. The reason was that we got
stuck to bounding $L^{(2)}_n$ rather than $R^{(2)}_k$, gradually
losing hope in an elegant solution.

Two explanations regarding our framework are in place.  First, the
Hilberg exponent of a sequence $(a_k)_{k\in\mathbb{N}}$ is defined as
\begin{align}
  \label{Exponent}
  \hilberg_{k\rightarrow\infty} a_k
  :=\limsup_{k\rightarrow\infty} \max\klam{0,\frac{\log
  a_k}{\log k}}
  .
\end{align}
It measures the rate of asymptotic power-law growth,
$\hilberg_{k\rightarrow\infty} k^\beta=\max\klam{0,\beta}$. Some
properties of the Hilberg exponent are summarized in \citep[Section
8.1]{Debowski21}. We write succinctly
\begin{align}
  a_k\lesssim b_k\iff
  \hilberg_{k\rightarrow\infty} a_k\le
  \hilberg_{k\rightarrow\infty} b_k,
  \\
  a_k\sim b_k\iff
  \hilberg_{k\rightarrow\infty} a_k=
  \hilberg_{k\rightarrow\infty} b_k.
\end{align}
The Hilberg exponent is useful for characterizing the power-law growth
of mutual information, i.e., Hilberg's law \citep[Chapter
8]{Debowski21}. Hence comes its name.

Second, for a parameter $\gamma\in(0,1)\cup(1,\infty)$, the
R\'enyi-Arimoto conditional entropy of a discrete random variable $X$
given a variable $Y$ is defined as
\begin{align}
  \label{RenyiArimoto}
  H_\gamma(X|Y):=
  -\frac{\gamma}{\gamma-1}\log
  \mean\kwad{\sum_x p(X=x|Y)^{\gamma}}^{1/\gamma}
  ,
\end{align}
whereas for $\gamma\in\klam{0,1,\infty}$, we define
\begin{align}
  H_\gamma(X|Y):=\lim_{\delta\rightarrow\gamma} H_\delta(X|Y).
\end{align}
We also put $H_\gamma(X):=H_\gamma(X|1)$. The special cases are
called:
\begin{itemize}
\item the Hartley entropy
$H_0(X|Y)= \log \esssup\card\klam{x: P(X=x|Y)>0}$,
\item  the Shannon
entropy $H_1(X|Y)=\mean\kwad{-\log P(X|Y)}$, 
\item the collision entropy
$H_2(X|Y)=-2\log \mean\kwad{\sum_{x\in\mathbb{X}} P(X=x|Y)^2}^{1/2}$,
\item the min-entropy $H_\infty(X|Y)=-\log \mean\max_{x} P(X=x|Y)$.
\end{itemize}
We have $H_\gamma(X)\ge H_\delta(X)$ for $\gamma<\delta$, whereas
$H_\gamma(X)\le \frac{\gamma}{\gamma-1}H_\infty(X)$ for $\gamma>1$.
Thus all R\'enyi-Arimoto entropies with $\gamma>1$ are of the same
magnitude.  Moreover, we have a generalized chain rule
\begin{align}
  \label{Chain}
  H_\gamma(X|Y,Z)\le H_\gamma(X|Y)\le H_\gamma(U,X|Y)\le
  H_0(U|Y)+H_\gamma(X|Y,U).
\end{align}
An overview of the R\'enyi-Arimoto entropies can be found in
\citep{FehrBerens14} and also in \citep[Section 9.1]{Debowski21}.

Let us come back to the recurrence and repetition times.  In the body of
the paper, we will prove the following Theorems \ref{theoMainOne} and
\ref{theoMainTwo}, which are analogous but a noticeable novelty can be
attributed only to Theorem \ref{theoMainTwo}.
\begin{theorem}
  \label{theoMainOne}
  Consider a stationary ergodic process $(X_i)_{i\in\mathbb{Z}}$ over
  a finite alphabet. We have a sandwich bound
  \begin{align}
    hk\sim
    \log P(X_1^k|X_{k+1}^\infty)\lesssim
    \log R^{(1)}_k\lesssim
    -\log P(X_1^k)\lesssim
    H_1(X_1^k)
    \text{ a.s.},
  \end{align}
  where the right-most inequality $\lesssim$ becomes the equivalence
  $\sim$ for
  \begin{align}
    \limsup_{k\to\infty} \frac{\sqrt{\var\log P(X_1^k)}}{H_1(X_1^k)}<1.
  \end{align}
\end{theorem}
\emph{Remark:} Theorem \ref{theoMainOne} applies the same sandwich
bound by \citet{Kontoyiannis98} as the simple proof of the WZ/OW
theorem (\ref{WZOW}) presented in \citep[Theorem 9.12]{Debowski21}.
We also recall that $h_1=\mean\kwad{-\log P(X_1|X_2^\infty)}$ for a
stationary process. Consequently if $h_1=0$ then
$P(X_1^k|X_{k+1}^\infty)=1$ almost surely. Quantity $\sqrt{\var\log
  P(X_1^k)}$ is called varentropy \citep{KontoyiannisVerdu13}.

\begin{theorem}
  \label{theoMainTwo}
  Consider a stationary ergodic process $(X_i)_{i\in\mathbb{Z}}$ over
  a countable alphabet. Define the context length
  \begin{align}
    \label{ContextLength}
    I_k:=\min\klam{i:\log i\ge H_\infty(X_1^k|X_{k+1}^{k+i})}.
  \end{align}
  We have a sandwich bound 
  \begin{align}
    \label{MainTwo}
    H_\infty(X_1^k|X_{k+1}^{k+I_k})\lesssim
    \log I_k\lesssim
    \log R^{(2)}_k\lesssim
    H_\infty(X_1^k)
    \text{ a.s.}
    ,
  \end{align}
  where the left-most inequality $\lesssim$ becomes the equivalence
  $\sim$ for a finite alphabet.
\end{theorem}
\emph{Remark:} We may substitute
$I_k\to\min\klam{i:\frac{\gamma}{\gamma-1}\log i\ge
  H_\gamma(X_1^k|X_{k+1}^{k+i})}$ and $H_\infty\to H_\gamma$ for any
$\gamma>1$ in the sandwich bound (\ref{MainTwo}).  In
\citep{Debowski18b,Debowski21}, a trivial bound in terms of Shannon
entropy $H_1(X_1^k)$ stemming from inequality $L^{(2)}_n\ge L^{(1)}_n$
was discussed instead of the right-most bound in (\ref{MainTwo}). By
contrast, the upper bound in (\ref{MainTwo}) in terms of min-entropy
$H_\infty(X_1^k)$, which is a brand new idea, follows by inequality
(\ref{Minimum}) and the generalized Kac lemma for successive
recurrence times by \citet{ChenMoy59}. As for the lower bound in
(\ref{MainTwo}), context length $I_k$ can be much shorter than the
pessimistic context length $D^k$ for a $D$-ary alphabet that was
proved sufficient in \citep[Theorem 8]{Debowski18b}. Moreover,
bounding $R^{(2)}_k$ rather than $L^{(2)}_n$ leads to more elegant
statements than produced in \citep{Debowski18b,Debowski21}.

The organization of the note is as follows. In Section
\ref{secRecurrence}, we recall the known sandwich bound for the
recurrence time. In Section \ref{secRepetition}, we prove a new
sandwich bound for the repetition time. Concluding, in Section
\ref{secFinal}, we specialize these results to the Hilberg exponents.

% Gaps to be filled:

% -- How much can different entropies differ?, Szpankowski's results
% on $L^{(\gamma)}_n$ for $\gamma\ge 3$ suggest that entropy rates
% $h_\gamma$ are equal.


\section{Sandwiching recurrence times}
\label{secRecurrence}

There is a celebrated lemma by \citeauthor{Kac47}. Namely, the
conditional expectation of the recurrence time is the inverse
probability of the initial string.
\begin{lemma}[\citep{Kac47}]
  \label{theoKac}
  For a stationary ergodic process $(X_i)_{i\in\mathbb{Z}}$ over a
  countable alphabet,
  \begin{align}
    \mean\okra{R^{(1)}_k\middle|X_1^k}
    =
    \frac{1}{P(X_1^k)}.
  \end{align}
\end{lemma}
There is also a converse result by \citeauthor{Kontoyiannis98}.
\begin{lemma}[\citep{Kontoyiannis98}]
  \label{theoKontoyiannis}
  For a process $(X_i)_{i\in\mathbb{Z}}$ over a $D$-ary alphabet,
  \begin{align}
    \mean\okra{\frac{1}{R^{(1)}_k
    P(X_1^k|X_2^\infty)}\middle|X_{k+1}^\infty}
    \le
    1+\log D^k.
  \end{align}
\end{lemma}
The origin of this clever but simple inequality comes from an
observation that for a fixed sequence $x_{k+1}^\infty$ and any number
$i\ge 1$, there is at most one string $x_1^k$ such that $R^{(1)}_k=i$
for $X_1^\infty=x_1^\infty$.

The following sandwich bound appears in the simple proof of the WZ/OW
theorem in \citep[Theorem 9.12]{Debowski21}. We recall the derivation
for the analogy with developments of Section \ref{secRepetition}.
\begin{proposition}
  \label{theoSandwichRecurrence}
  Consider a stationary ergodic process $(X_i)_{i\in\mathbb{Z}}$ over
  a $D$-ary alphabet. Assume that $\sum_{k=1}^\infty \rho_k<\infty$
  for some $\rho_k>0$.  Almost surely, for sufficiently large $k$ we have
\begin{align}
  \label{SandwichRecurrence}
  -\log P(X_1^k|X_{k+1}^\infty)
  +\log \rho_k-\log k
  <
  \log R^{(1)}_k
  <
  -\log P(X_1^k)
  -\log \rho_k
  .
\end{align}
\end{proposition}
\begin{proof}
  As \citet{Kontoyiannis98} has shown, by applying the Markov
  inequality to Lemmas \ref{theoKac} and \ref{theoKontoyiannis}, we
  obtain two complementary bounds
\begin{align}
  \label{ProbRecurrenceI}
  P\okra{R^{(1)}_k\ge \frac{C}{P(X_1^k)}}&\le C^{-1}
  ,
  \\
  \label{ProbRecurrenceII}
  P\okra{R^{(1)}_k\le \frac{C}{P(X_1^k|X_{k+1}^\infty)}}&\le C(1+k\log D)
  .
\end{align}
Setting $C=\rho_k^{-1}$ in (\ref{ProbRecurrenceI}) and $C=k^{-1}\rho_k$ in
(\ref{ProbRecurrenceII}) we obtain (\ref{SandwichRecurrence}) by the
Borel-Cantelli lemma.
\end{proof}

The WZ/OW theorem (\ref{WZOW}) follows by Proposition
\ref{theoSandwichRecurrence} combined with the
Shannon-McMillan-Breiman (SMB) theorem \citep{Shannon48, Breiman57,
  Chung61, Barron85, AlgoetCover88}, which asserts
\begin{align}
  \label{SMB}
  \lim_{k\to\infty}
  \frac{1}{k}\okra{-\log P(X_1^k)}=
  \lim_{k\to\infty}
  \frac{1}{k}\okra{-\log P(X_1^k|X_{k+1}^\infty)}=
  h_1 \text{ a.s.}
\end{align}

\section{Sandwiching repetition times}
\label{secRepetition}

There is a result by \citeauthor{ChenMoy59} that generalizes the Kac
lemma to successive recurrence times. In our notation, it can be
stated in the following way.
\begin{lemma}[\citep{ChenMoy59}]
  \label{theoChenMoy}
  For a stationary process $(X_i)_{i\in\mathbb{Z}}$, fix a string
  $x_1^k$ such that $P(X_1^k=x_1^k)>0$. Define the random set
  $\mathcal{T}:=\klam{j\in\mathbb{Z}: X_{j}^{j+k-1}=x_1^k}$.  Let the
  sequence $(T_r)_{r\in\mathbb{Z}}$ be defined by conditions:
  $T_r\in\mathcal{T}$ and $T_r>T_{r-1}$ for $r\in\mathbb{Z}$, whereas
  $T_1\ge 1$ and $T_0\le 0$. The successive recurrence times are
  defined as
  \begin{align}
    \label{Successive}
    W_r:=T_r-T_{r-1}.
  \end{align}
  It is claimed that for a stationary process $(X_i)_{i\in\mathbb{Z}}$
  with respect to measure $P$, process $(W_r)_{r\in\mathbb{Z}}$ is
  stationary with respect to measure $P(\cdot|T_0=0)$. Moreover, if
  additionally process $(X_i)_{i\in\mathbb{Z}}$ is ergodic then for
  $r\in\mathbb{Z}$, we have
  \begin{align}
    \mean\okra{W_r\middle|T_0}
    =
    \frac{1}{P(X_1^k=x_1^k)}.
  \end{align}
\end{lemma}

Let us derive a strong upper bound for the repetition time.  The
following proposition is a new simple idea. We have sought for this
sort of bound in terms of the maximal repetition length intermittently
for a long time. In our perception, this involved lower bounding the
probability of a union for strongly dependent processes. We tried to
use the Chung-Erd\H{o}s inequality \citep{ChungErdos1952} or other
second-moment methods like \citet{Szpankowski93,Szpankowski93a} did,
which led to complicated calculations and assumptions difficult to
satisfy. We have overlooked that the solution may be so simple when we
reformulate the problem as bounding the repetition time.
\begin{proposition}
  \label{theoSandwichRepetitionOne}
  Consider a stationary ergodic process
  $(X_i)_{i\in\mathbb{Z}}$ over a countable alphabet. Assume that
  $\sum_{k=1}^\infty \rho_k<\infty$ for some $\rho_k>0$.  Almost
  surely, for sufficiently large $k$ we have
\begin{align}
  \label{SandwichRepetitionOne}
  \log R^{(2)}_k< H_\infty(X_1^k)-\log\rho_k
  .
\end{align}
\end{proposition}
\begin{proof}
  Consider an arbitrary string $x_1^k$ such that $P(X_1^k=x_1^k)>0$
  and the corresponding successive recurrence times
  (\ref{Successive}).  Specializing inequality (\ref{Minimum}), we
  obtain
  \begin{align}
    R^{(2)}_k\le W_1+W_2,
  \end{align}
  since by the time we observe string $x_1^k$ twice, we must have seen
  some repetition of length $k$, certainly.  Fix $x_1^k$ as the most
  probable string of length $k$.  In view of Lemma \ref{theoChenMoy},
  we may bound
  \begin{align}
    P(R^{(2)}_k\ge n)
    &\le \frac{\mean R^{(2)}_k}{n}
      \le \frac{\mean \mean\okra{W_1+W_2|T_0}}{n}
      \nonumber\\
    &= \frac{2}{nP(X_1^k=x_1^k)}
      = \frac{2e^{-H_\infty(X_1^k)}}{n},
  \end{align}
  since $e^{-H_\infty(X)}=\max_{x} P(X=x)$.  Define
  $n_k:=e^{H_\infty(X_1^k)}\rho_k^{-1}$. Then we have
  $\sum_{k=1}^\infty P(R^{(2)}_k\ge n_k)<\infty$. Hence inequality
  (\ref{SandwichRepetitionOne}) is obtained by the Borel-Cantelli
  lemma.
\end{proof}

Now we will derive a lower bound for the repetition time.  This
proposition refines \citep[Theorem 9.25]{Debowski21} by using the
weighted conditional entropy rather than the min-entropy conditioned
on the infinite future.
\begin{proposition}
  \label{theoSandwichRepetitionTwo}
  Consider a stationary process $(X_i)_{i\in\mathbb{Z}}$ over a
  countable alphabet. Assume that $\sum_{k=1}^\infty \rho_k<\infty$
  for some $\rho_k>0$. Define the weighted conditional entropy
  \begin{align}
    H_\bullet(X_1^k|X_{k+1}^\infty):=-\log \okra{
    \sum_{i=1}^\infty\frac{e^{-H_\infty(X_1^k|X_{k+1}^{k+i})}}{i(i+1)}}.
  \end{align}
  Almost surely, for sufficiently large $k$ we have
  \begin{align}
    \label{SandwichRepetitionTwo}
    \log R^{(2)}_k> \frac{1}{3}H_\bullet(X_1^k|X_{k+1}^\infty)
    +\frac{1}{3}\log\rho_k
    .
  \end{align}
\end{proposition}
\emph{Remark:} We note that
$H_\bullet(X_1^k|X_{k+1}^\infty)\ge
H_\infty(X_1^k|X_{k+1}^\infty)$. Hence Proposition
\ref{theoSandwichRepetitionTwo} is stronger than \citep[Theorem
9.25]{Debowski21}.
\begin{proof}
  For $i<j$, we obtain 
\begin{align}
  P(X_{i+1}^{i+k}=X_{j+1}^{j+k})
  &=
    \mean
    P(X_{i+1}^{i+k}=X_{j+1}^{j+k}|X_{i+k+1}^{j+k})
    \nonumber\\
  &\le
    \mean \max_{x_1^k} P(X_{i+1}^{i+k}=x_1^k|X_{i+k+1}^{j+k})
    \nonumber\\
  &=
    e^{-H_\infty(X_1^k|X_{k+1}^{k+j-i})}
    .
\end{align}
Observe that $3-2x\le x^{-2}$, so $n-x\le 4n^3x^{-2}/27$.
%$n-x=n/3(3-2(3x/2n))\le n/3(3x/2n)^{-2}=4n^3/27x^2$.
Hence, by the union bound,
\begin{align}
  P(R^{(2)}_k\le n)
  &= P\okra{X_{i+1}^{i+k}=X_{j+1}^{j+k} \text{ for some }
    0\le i<j\le n}
    \nonumber\\
  &\le \sum_{0\le i<j\le n} P(X_{i+1}^{i+k}=X_{j+1}^{j+k})
    \le \sum_{0\le i<j\le n}  e^{-H_\infty(X_1^k|X_{k+1}^{k+j-i})}
    \nonumber\\
  &= \sum_{i=1}^{n-1}
    \frac{n-i}{e^{H_\infty(X_1^k|X_{k+1}^{k+i})}}
    \le \frac{8n^3}{27}
    \sum_{i=1}^\infty\frac{e^{-H_\infty(X_1^k|X_{k+1}^{k+i})}}{i(i+1)}
    \nonumber\\
  &\le \frac{8n^3}{27} e^{-H_\bullet(X_1^k|X_{k+1}^\infty)}.
\end{align}
Put $n_k:=\sqrt[3]{\rho_k e^{H_\bullet(X_1^k|X_{k+1}^\infty)}}$. Then
$\sum_{k=1}^\infty P(R^{(2)}_k\le n_k)<\infty$.  Consequently,
inequality (\ref{SandwichRepetitionTwo}) follows by the
Borel-Cantelli lemma.
\end{proof}

The last proposition in this section is also new although it recycles
some techniques used in universal coding such as \citep[Theorem
6.11]{Debowski21}. The goal is to show that the bound in terms of the
weighted conditional entropy is more precise than the bound proved in
\citep[Theorem 8]{Debowski18b} since the length of conditioning can be
shortened to $I_k$ rather than $D^k$.
\begin{proposition}
  \label{theoSandwichContext}
  Consider a stationary process $(X_i)_{i\in\mathbb{Z}}$ over a
  countable alphabet.  For $I_k$ defined in (\ref{ContextLength}), we
  have
\begin{align}
  \label{SandwichContextOne}
  \log I_k-\log 2&\le H_\bullet(X_1^k|X_{k+1}^\infty)\le
  3\log I_k+\frac{1}{I_k}
  ,
  \\
  \label{SandwichContextThree}
  \log (I_k-1)-H_0(X_i)&\le H_\infty(X_1^k|X_{k+1}^{k+I_k})\le \log I_k
                              .
\end{align}
\end{proposition}
\emph{Remark:} We note that $I_k\le D^k$ and
$H_\infty(X_1^k|X_{k+1}^{k+I_k})\ge H_\infty(X_1^k|X_{k+1}^{k+D^k})$
for a $D$-ary alphabet. Hence Propositions
\ref{theoSandwichRepetitionTwo} and \ref{theoSandwichContext} combined
are stronger than \citep[Theorem 8]{Debowski18b}.
\begin{proof}
First, we derive (\ref{SandwichContextOne}) from
\begin{align}
  H_\bullet(X_1^k|X_{k+1}^\infty)
  &\le
    H_\infty(X_1^k|X_{k+1}^{k+I_k})+\log I_k+\log (I_k+1)\le
    3\log I_k+\frac{1}{I_k},
  \\
  H_\bullet(X_1^k|X_{k+1}^\infty)
  &\ge
    -\log\okra{\okra{1-\frac{1}{I_k}}e^{-H_\infty(X_1^k|X_{k+1}^{k+I_k-1})}
    +\frac{1}{I_k}}
    \nonumber\\
  &\ge -\log\okra{\okra{1-\frac{1}{I_k}}\frac{1}{I_k-1}+\frac{1}{I_k}}
    =\log I_k-\log 2.
\end{align}
Moreover, for a stationary process,
\begin{align}
  0
  &\le
    H_\infty(X_1^k|X_{k+1}^{k+I_k-1})-H_\infty(X_1^k|X_{k+1}^{k+I_k})\le
    H_0(X_i)
\end{align}
since by the chain rule (\ref{Chain}), we have
\begin{align}
  H_\infty(Y,X|Z,V)\le H_\infty(Y,X|Z)
  &\le H_0(X|Z)+H_\infty(Y|X,Z)
    \nonumber\\
  &\le H_0(X)+H_\infty(U,Y|X,Z). 
\end{align}
Hence we obtain (\ref{SandwichContextThree}).
\end{proof}

\section{Final touch}
\label{secFinal}

Let us recall some properties of Hilberg exponents.  The probabilistic
results of \citep[Section 8.1]{Debowski21} can be summarized as
follows.
\begin{definition}
  Let $T$ be an automorphism of the probability space.  A sequence of
  random variables $(J_k)_{k\in\mathbb{N}}$ is called $T$-increasing
  if
  \begin{align}
    J_{k+1}\ge J_k\circ T,J_k\ge 0.
  \end{align}
\end{definition}
Let $\median X:=\sup \klam{r: P(X<r)\le \frac{1}{2}}$ denote the
median of a variable $X$.
\begin{lemma}
  \label{theoHilberg}
  For $T$ being a stationary ergodic automorphism and a $T$-increasing
  sequence $(J_k)_{k\in\mathbb{N}}$, we have
  \begin{align}
    \hilberg_{k\rightarrow\infty} \median J_k
    \le \hilberg_{k\rightarrow\infty} J_k 
    \le \hilberg_{k\rightarrow\infty} \mean J_k
    \text{ a.s.}, 
  \end{align}
  where random variable $\hilberg_{k\rightarrow\infty} J_k$ is almost
  surely constant. Moreover, the inequalities become the equalities if
  \begin{align}
    \limsup_{k\rightarrow\infty} \frac{\sqrt{\var J_k}}{\mean J_k}<1.
  \end{align}
\end{lemma}
\begin{proof}
  All claims but the last one follow by \citep[Theorems 8.3--8.5]{Debowski21}.
  The last claim is a slight strengthening of \citep[Theorem
  8.6]{Debowski21} and in a similar fashion it follows by the
  median-mean inequality $\abs{\median X-\mean X}\le \sqrt{\var X}$,
  see e.g.\ \citep[Theorem 3.17]{Debowski21}.
\end{proof}

We also make this simple observation.
\begin{proposition}
  \label{theoIncreasing}
  Consider the shift automorphism $X_i\circ T=X_{i+1}$. Then sequences
  $(L^{(\gamma)}_n)_{n\in\mathbb{N}}$,
  $(R^{(\gamma)}_k)_{k\in\mathbb{N}}$,
  $(-\log P(X_1^k))_{k\in\mathbb{N}}$, and
  $(-\log P(X_1^k|X_{k+1}^\infty))_{k\in\mathbb{N}}$ are
  $T$-increasing.
\end{proposition}
Theorem \ref{theoMainOne} follows by Propositions
\ref{theoSandwichRecurrence} and \ref{theoIncreasing} and Lemma
\ref{theoHilberg}, whereas Theorem \ref{theoMainTwo} follows by
Propositions
\ref{theoSandwichRepetitionOne}--\ref{theoSandwichContext}. Quod erat
demonstrandum.

% \bibliographystyle{IEEEtran}
\bibliographystyle{abbrvnat}
%\bibliographystyle{elsart-harv}

\bibliography{0-journals-abbrv,0-publishers-abbrv,ai,mine,tcs,ql,nlp,books}

\end{document}

% The subword complexity \citep{Ferenczi99, DeLuca99,
%   JansonLonardiSzpankowski04, GheorghiciucWard07, Ivanko08} is the
% number of distinct substrings of a given length in a given string,
% \begin{align}
%  f(k|x_1^n):=\card\klam{x_{i+1}^{i+k}:0\le i\le n-k}. 
% \end{align}
% We have two simple facts.
% \begin{lemma}[\mbox{\citep[Theorem 9.15]{Debowski21}}]
%   \label{theoSCRecurrence}
%   For any stationary process $(X_i)_{i\in\mathbb{Z}}$, 
% \begin{align}
%   \mean f(k|X_1^{n+k-1})=\mean \min\klam{R^{(1)}_k,n}.
% \end{align}
% \end{lemma}
% \begin{lemma}[\mbox{\citep[Theorem 9.21]{Debowski21}}]
%   \label{theoSCRepetition}
%   For any process $(X_i)_{i\in\mathbb{Z}}$, 
%   \begin{align}
%     P(L^{(2)}_{n+k-1}<k)\le \frac{\mean f(k|X_1^{n+k-1})}{n}.
%   \end{align}
% \end{lemma}
% Let us observe that Lemma \ref{theoSCRepetition} can be alternatively
% proved using Lemma \ref{theoSCRecurrence} and the Markov inequality,
% \begin{align}
%   P(L^{(2)}_{n+k-1}<k)
%   &\le P(R^{(2)}_k\ge n)
%   =P(R^{(1)}_k\ge n)
%     =P\okra{\min\klam{R^{(1)}_k,n}\ge n}
%     \nonumber\\
%   &\le \frac{\mean\min\klam{R^{(1)}_k,n}}{n}
%   = \frac{\mean f(k|X_1^{n+k-1})}{n}.
% \end{align}
% Thus it cannot lead to a non-trivial upper bound for $R^{(2)}_k$.
 
% Now, we will derive an upper bound for the repetition time.  The
% following proposition is a new idea.
% \begin{proposition}
%   \label{theoSandwichRepetitionOne}
%   Consider a stationary ergodic process $(X_i)_{i\in\mathbb{Z}}$.  Let
%   us denote
% \begin{align}
%   E_k&:=\mean P(X_1^k)=e^{-H_2(X_1^k)},
%   \\
%   V_k&:=\var P(X_1^k)=e^{-2H_3(X_1^k)}-e^{-2H_2(X_1^k)}.
% \end{align}
% Assume that $V_k/E_k^2\le \rho_k$ where
% $\sum_{k=1}^\infty \rho_k\log k<\infty$ for some $\rho_k>0$. Almost
% surely, for sufficiently large $k$ we have
% \begin{align}
%   \label{SandwichRepetitionOne}
%   \log R^{(2)}_k< H_2(X_1^k)-\log\rho_k
%   .
% \end{align}
% \end{proposition}
% \begin{proof}
%   By Lemmas \ref{theoSCRecurrence} and \ref{theoKac}, we obtain
% \begin{align}
%   \mean f(k|X_1^{n+k-1})
%   &
%     =\mean \min\klam{R^{(1)}_k,n}
%     %\nonumber\\
%   %&
%     =
%     \mean\mean\okra{\min\klam{R^{(1)}_k,n}\middle|X_1^k}
%     \nonumber\\
%  &\le
%    \mean\min\klam{\mean(R^{(1)}_k|X_1^k),n}
%    %\nonumber\\
%   %&
%     =
%     \mean\min\klam{\frac{1}{P(X_1^k)},n}
%     \nonumber\\
%   &=
%     \int_0^{n} P\okra{\frac{1}{P(X_1^k)}\ge x} dx
%     %\nonumber\\
%   %&
%     =
%      n-\int_0^{n} P\okra{P(X_1^k)>x^{-1}} dx
%      .
% \end{align}
% Recall the Paley-Zygmund inequality, see e.g.\ \citep[Theorem
% 3.16]{Debowski21},
% \begin{align}
%   P(Y> \epsilon)\ge
%   \frac{(\mean Y-\epsilon)_+^2}{\var Y+(\mean Y-\epsilon)_+^2}
%   .
% \end{align}
% For $n\ge 1/E_k$, we obtain
% \begin{align}
%   \mean f(k|X_1^{n+k-1})
%   &\le
%     n-\int_0^{n}\frac{(\mean P(X_1^k)-x^{-1})^2_+
%     }{\var P(X_1^k)+(\mean P(X_1^k)-x^{-1})^2_+} dx
%     \nonumber\\
%   &=
%     n-\int_{1/E_k}^{n}\frac{(E_kx-1)^2}{V_kx^2+(E_kx-1)^2}dx.
%     \nonumber\\
%   &\le
%     n-\int_{1/E_k}^{n} \frac{E_k^2x^2-2E_kx}{V_kx^2+E_k^2x^2} dx
%      \nonumber\\
%   &=
%     n -\frac{E_k^2}{V_k+E_k^2}\okra{n-\frac{1}{E_k}}
%     +\frac{2E_k}{V_k+E_k^2}\log E_kn
%     \nonumber\\
%   &=
%     \frac{V_kn+E_k(1+2\log E_kn)}{V_k+E_k^2}.
% \end{align}
% Applying Lemma \ref{theoSCRepetition}, we may rewrite
% \begin{align}
%   P(R^{(2)}_k\ge n)=P(L^{(2)}_{n+k-1}<k)
%   &\le
%     \frac{\mean f(k|X_1^{n+k-1})}{n}
%     \nonumber\\
%   &\le
%   \frac{V_k}{V_k+E_k^2}+\frac{E_k^2}{V_k+E_k^2}\frac{1+2\log
%     E_kn}{E_kn}
%     .
% \end{align}
% Define $n_k:=E_k^{-1}\rho_k^{-1}$. Then
% \begin{align}
%   P(R^{(2)}_k\ge n_k)\le
%   \frac{\rho_k}{1+\rho_k}+\frac{\rho_k(1-2\log
%  \rho_k)}{1+\rho_k}.
% \end{align}
% In particular, $\sum_{k=1}^\infty P(R^{(2)}_k\ge n_k)<\infty$. Hence
% inequality (\ref{SandwichRepetitionOne}) is obtained by the
% Borel-Cantelli lemma.
% \end{proof}


% We can check the indefinite integral
% \begin{align}
%   &\int \frac{(Ex-1)^2}{Qx^2-2Ex+1}dx
%     =
%   \nonumber\\
%   &\quad\frac{E^2x}{Q}-\frac{VE}{Q^2}\log\okra{Qx^2-2Ex+1}
%     +\frac{\sqrt{V}(V-E^2)}{Q^2}\tan^{-1}\okra{\frac{Qx-E}{\sqrt{V}}}
%     .
% \end{align}
% Hence
% \begin{align}
%   \int_{1/E}^{n_k}\frac{(Ex-1)^2}{Qx^2-2Ex+1}dx
%   &\ge
%     \frac{E^2n_k}{Q}
%     -\frac{VE}{Q^2}\log\okra{Qn_k^2-2En_k+1}+C,
%   \\
%   C
%   &:=
%     -\frac{E}{Q}
%     +\frac{VE}{Q^2}\log\okra{\frac{V}{E^2}}
%     -\frac{\pi\sqrt{V}(V-E^2)}{Q^2}
%     .  
% \end{align}

% Denote $A_{ij}:=X_{i+1}^{i+k}=X_{j+1}^{j+k}$,
% $E:=\mean\max_{x_1^k} P(X_{i+1}^{i+k}=x_1^k|X_{i+k+1}^{\infty})$,
% $S:=\sum_{0\le i<j\le n} P(A_{ij})$, $N:=n(n-1)/2$,
% $\pi_{ijlm}\propto P(A_{ij}\cap A_{lm})$,
% $\psi_{ijlm}=\frac{P(A_{ij}\cap A_{lm})}{P(A_{ij})P(A_{lm})}$.  By the
% Chung-Erd\H{o}s inequality,
% \begin{align}
%   P(R^{(2)}_k\le n)
%   &= P\okra{\bigcup_{0\le i<j\le n} A_{ij}}
%     \ge 
%     \frac{
%     \displaystyle\sum_{0\le i<j\le n} \sum_{0\le l<m\le n} 
%     P(A_{ij})P(A_{lm})
%     }{
%     \displaystyle\sum_{0\le i<j\le n} \sum_{0\le l<m\le n}
%     P(A_{ij}\cap A_{lm})
%     }
%     \nonumber\\
%   &=
%     \sum_{ijlm}
%     \pi_{ijlm}\okra{\psi_{ijlm}}^{-1}
%     \ge
%     \okra{\sum_{ijlm}
%     \pi_{ijlm}\psi_{ijlm}}^{-1}
%   %   \nonumber\\
%   % &\ge \okra{\max_{0\le i<j\le n}\max_{0\le i<j\le n}
%   %   \frac{P(A_{ij}\cap A_{lm})}{P(A_{ij})P(A_{lm})}}^{-1}
%     \nonumber\\
%   &\ge
%     \frac{S^2}{S(1+(N-1)E)}
%     \ge
%     \frac{S}{1+(N-1)E}
%     \nonumber\\  
%   &=\frac{
%     S_{i<j}S_{l<m}
%     }{
%     S_{i=l<j=m}
%     +2S_{i=l<j<m}+2S_{l<m=i<j}+2S_{i<l<j=m}
%     +2S_{l<i<j<m}+2S_{i<l<j<m}+2S_{i<j<l<m}
%     }
%    \nonumber\\
%   &=\frac{
%     S_{i<j}^2
%     }{
%     S_{i<j}
%     +6S_{i=l<j<m}
%     +2S_{l<i<j<m}+2S_{i<l<j<m}+2S_{i<j<l<m}
%     }
%   .  
% \end{align}

% \begin{lemma}[\mbox{\citep[Theorem 8.3]{Debowski21}}]
%   \label{theoHilbergConstant}
%   Consider a sequence of real random variables
%   $(J_k)_{k\in\mathbb{N}}$.  If $T$ is a stationary ergodic
%   transformation of the probability space and $J_{k+1}\ge J_k\circ T$
%   for all $k\in\mathbb{N}$ then random variable
%   $\hilberg_{k\rightarrow\infty} J_k$ is almost surely constant.
% \end{lemma}

% \begin{lemma}[\mbox{\citep[Theorem 8.4]{Debowski21}}]
%   \label{theoHilbergExpectation}
%   Consider a sequence of real random variables
%   $(J_k)_{k\in\mathbb{N}}$.  If $0\le J_k\le J_{k+1}$ for all
%   $k\in\mathbb{N}$ then
% \begin{align}
%   \label{HilbergExpectation}
%   \esssup\hilberg_{k\rightarrow\infty} J_k \le
%   \hilberg_{k\rightarrow\infty} \mean J_k 
%   .
% \end{align}
% \end{lemma}

% \begin{lemma}[\mbox{\citep[Theorem
%     8.5]{Debowski21}}]
%   \label{theoHilbergMedian}
%   Consider a sequence of real random variables
%   $(J_k)_{k\in\mathbb{N}}$ such that $J_k\ge 0$ for all
%   $k\in\mathbb{N}$.  We have
%   \begin{align}
%   \label{HilbergMedian}
%   \hilberg_{k\rightarrow\infty} \median J_k
%   \le \esssup \hilberg_{k\rightarrow\infty} J_k 
%   .
% \end{align}
% \end{lemma}


%%% Local Variables: 
%%% mode: latex
%%% mode: flyspell
%%% TeX-master: t
%%% TeX-PDF-mode: t
%%% ispell-local-dictionary: "american"
%%% End: 

%  LocalWords:  differentiable
